\section{Physical coordinates and barycentric sweeps}
\label{sec:coordinate-realizations}

This appendix translates the caps and finite feet of
Appendices~\ref{sec:cap-choices}--\ref{sec:finite-footprints} into other
Euclidean realizations of a regular tetrahedron. Each translation must
track both the physical area factor and the effective tilt: rescaling an
angular parameter alone does not determine either. A single dictionary
handles the scaled models. We then give a barycentric realization, where
a direct overlap proof and a geometric enclosure also explain positivity
of the radial parameters and coverage of boundary points.

\subsection{A scaled coordinate model}

Let $b>0$ and $H_b=\sqrt2 b$. In physical coordinates
$(x,y,H_bt)$, define the regular edge-$2b$ tetrahedron
\[
 K_b=\{(x,y,H_bt):0\le t\le1,\quad
            |x|\le bt,\quad |y|\le b(1-t)\}.
\]
Its paired line families, with angular variable $s\in[-b,b]$, are
\begin{align*}
 &(x,\rho s-\eta x(s^2+b^2)/2,
                  H_b(1-\rho+\eta xs)),\\
 &(\rho s+\eta y(s^2+b^2)/2,y,
                  H_b(\rho+\eta ys)).
\end{align*}
Their free coordinate satisfies $|x|\le b$ or $|y|\le b$, and
$0\le\rho\le T_\eta(s)$. The transformation to the earlier physical
model is
\begin{equation}\label{eq:halfheight-dictionary}
 X=2x/b,\quad Y=2y/b,\quad Z=2t-1,
 \qquad r=s/b,\quad R=2\rho,\quad \eps=\eta b^2.
\end{equation}
The canonical physical point is $(X,Y,\sqrt2 Z)$. This map is a
translation followed by a Euclidean similarity of ratio $2/b$.
Substituting in the first displayed line gives
$(X,Rr-\eps XG(r),1-R+\eps Xr)$ in the $(X,Y,Z)$ coordinates;
substitution in the second gives the lower line. This checks the two
families and their radial bounds, not only the tetrahedron.
The common upper bound is $R_\eps(r)=2T_\eta(br)$.

For example, cap (L) and cap (S), respectively, become
\begin{align}
 T_\eta(s)&=\frac12+\eta^2 A_{b,\beta}(s),\label{eq:scaled-linear-cap}\\
 T_\eta(s)&=\sqrt{\frac14+\eta^2 A_{b,\beta}(s)},\label{eq:scaled-square-cap}\\
 A_{b,\beta}(s)&=\frac{s^4+b^2s^2}{4}+\frac{b^4\beta}{2}.
 \nonumber
\end{align}
The exact two-family area integral in these coordinates is
\begin{equation}\label{eq:scaled-cost-identity}
 H_b\int_{-b}^b
 \frac{T_\eta(s)^2}{\sqrt{1+\eta^2((s^2+b^2)^2/4+H_b^2s^2)}}\,ds
 =\frac{bH_b}{2}Q_R(\eta b^2).
\end{equation}
Indeed $ds=b\,dr$, $T=R/2$, and the denominator becomes
$\sqrt{1+\eps^2D_0(r)}$. Both sides express the same weighted cost in physical area units. The multiplier $bH_b/2$ is half the minimum projection area
of $K_b$. Thus the coefficient in these coordinates, relative to this
reference area, is exactly $(2\beta-1/30)\eta^2b^4$.

The following examples use tetrahedra of edge lengths one, $\sqrt2$,
and two. Keeping the physical scale and angular scale together lets
each row determine both the finite bases and their cost coefficients.
The parameter $\tau$ denotes the tilt used in the indicated row.
\begin{center}
\begin{tabular}{llll}
$b$ & $\eta$ & $\beta$ & radial cap\\ \hline
$1$ & $\tau$ & $1/120$ & linear\\
$1/2$ & $\tau/2$ & $1/120$ & linear\\
$1/\sqrt2$ & $\tau$ & $1/125$ & linear or square root\\
$1/2$ & $\tau$ & $4/625$ & linear\\
$1/2$ & $2\tau$ & $1/120$ & square root\\
$1$ & $2\tau$ & $1/200$ & linear
\end{tabular}
\end{center}
For instance, the $b=1/\sqrt2$ row has additive polynomial
$(s^4+s^2/2)/4+1/1000$ and relative quadratic coefficient
$-13\tau^2/3000$. The $b=1/2$, $\eta=\tau$ row has additive
polynomial $s^4/4+s^2/16+1/5000$, and its limiting total area is
\[
 2H_b\left(\frac18-\frac{77}{480000}\tau^2+O(\tau^4)\right).
\]
For $b=1/2$, $\eta=\tau/2$, write $c=1/2$ and
$a(s)=c(1/8+s^2/2)$; the cap can equivalently be written
$c+\tau^2(c^2s^2a(s)+c^3/1920)$. Its limiting total area is
$H_b/4-(H_bc^4/960)\tau^2+O(\tau^4)$.
The monotone-polynomial calculation in
Appendix~\ref{sec:cap-choices} applies with precisely these constants.

Under \eqref{eq:halfheight-dictionary}, the generously padded
absolute-value footprint $A=B=\sqrt2+2\eps$ and the midpoint affine
footprint $A=B=4\eps$ remain explicit finite enclosures.
Equations~\eqref{eq:abs-foot-area} and \eqref{eq:affine-foot-area} compute
their finite areas; the sector lower bound in Theorem~\ref{thm:radial}
shows that actual swept projections have the same limiting total area.
A similarity preserves that limit, the possible differences between the
finite areas, convexity, and the number $2N$ of cylinders.

\paragraph{Centered coordinates.}
For a centered tetrahedron one instead uses
$|x|\le b(1+t)$, $|y|\le b(1-t)$ with physical height coordinate
$H_bt$. The similarity $X=x/b$, $Y=y/b$, $Z=t$ has ratio $1/b$.
Angular substitution $r=s/b$ again makes the effective tilt
$\eps=\eta b^2$; the radial parameter now stays unchanged.
Its two-family reference cost is $2\sqrt2 b^2$. For example, when
$b=1/\sqrt2$, cap (L) with $\beta=1/250$ becomes
\[
 R_\eta(s)=1+\eta^2\left((s^4+s^2/2)/2+1/1000\right).
\]
Half the total quadratic coefficient in these orthonormal coordinates
is $-b/120+2b/1000<0$. Fractional padding
$R_\eta(s)=1+\eta^2s^2(b^2+s^2)/2+\eta^{5/2}$ becomes cap (F)
with $c=4\sqrt2$; its limiting total cost is
$2b-b\eta^2/60+o(\eta^2)$. These formulas retain the factor of two
between one family and the full cover.

\paragraph{Nonorthonormal coordinates.}
One useful coordinate system has metric
$dt^2+(dp^2+dq^2)/2$ and tetrahedron
$|p|\le1+t$, $|q|\le1-t$. Sending it to the canonical physical point
$(p,q,\sqrt2 t)$ is a similarity of ratio $\sqrt2$, rather than an
isometry. Its physical intercept area multiplier is $1/\sqrt2$;
its limiting total cost is $\sqrt2 Q_R(\eps)$, and its minimum projection
area is $2\sqrt2$. Cubic padding $5\eps^3$ therefore gives
\[
 \sqrt2\left(1-\eps^2/30+10\eps^3+O(\eps^4)\right),
\]
as the exact limit for both the explicit feet and the actual swept
projections. Fractional padding $\eps^{5/2}$ gives
$\sqrt2(1-\eps^2/30+2\eps^{5/2}+O(\eps^4))$.
\paragraph{An orthonormal height coordinate.}
Consider the physical tetrahedron $|x|\le1-kz$, $|y|\le1+kz$ in
ordinary orthonormal coordinates, where $k=1/\sqrt2$. Its canonical
coordinates are $X=y$, $Y=x$, $t=kz$. For a positive physical tilt
$\tau$, the effective canonical tilt is $\eps=k\tau$.
Write the physical radial height cap as
\[
 R_{\mathrm{phys}}(r)=k^{-1}+k\tau^2H(r)+10\tau^3.
\]
Its canonical cap is $kR_{\mathrm{phys}}$, so $\eps=k\tau$ and
$k^2=1/2$ give canonical cubic constant $c=20$.
Select the angular roots $q$ for the canonical upper family and $p$ for
the canonical lower family as in Proposition~\ref{prop:coverage}.
If both selected radial parameters exceed their caps, then after the
coordinate substitution, \eqref{eq:cap-failure-bound} multiplied by
$\eps\tau$ gives
\[
 \tau(yq-xp)\le k\tau^2(H(p)+H(q))+9k^2\tau^3.
\]
The two physical cubic additions supply $20\tau^3$, and $9k^2<20$.
This preserves the explicit cubic comparison in physical, rather than
rescaled, tilt units.

\subsection{A barycentric realization and its boundary points}

In this subsection $k$ denotes a new positive tilt parameter, independent
of the fixed conversion constant used above. It will be chosen sufficiently
small.

Consider the nonnegative simplex $\Delta$ in the affine Euclidean space
\[
 \{(a,b,c,d)\in\mathbb R^4:a+b+c+d=1\}.
\]
Its edge length is $\sqrt2$. Set $s=a+b$, $q=a-b$, and $p=c-d$.
The induced squared length element is
$ds^2+(dq^2+dp^2)/2$. Hence
\begin{equation}\label{eq:barycentric-dictionary}
 (X,Y,Z)=(2q,2p,2s-1)
\end{equation}
gives a Euclidean similarity of ratio $2\sqrt2$ to the canonical
physical model, where the physical point is $(X,Y,\sqrt2 Z)$.
The minimum projection area of $\Delta$ is therefore $1/\sqrt2$.

Define, for $0\le t\le1$,
\[
 V(t)=(t^2,-(1-t)^2),\quad w(t)=2t-1,\quad
 F(v)=v^2+4v^4+1/1000,\quad
 B_k(t)=1/2+k^2F(w(t)/2).
\]
For the first sweep set
\begin{equation}\label{eq:barycentric-first-sweep}
 (a,b)=S(t,1-t)+kpV(t),\quad
 0\le S\le B_k(t),\quad |p|\le1,
\end{equation}
with $c,d$ determined by their sum and difference. For the second set
\begin{equation}\label{eq:barycentric-second-sweep}
 (c,d)=T(r,1-r)-kqV(r),\quad
 0\le T\le B_k(r),\quad |q|\le1.
\end{equation}
Under \eqref{eq:barycentric-dictionary}, the first family is the
canonical lower family and the second is the upper family, with angular
parameters $2t-1$ and $2r-1$, radial parameters $2S,2T$, and free
coordinates $2p,2q$. Their cap is precisely
$1+k^2(H+1/500)$. This already gives coverage for small $k$.
The following direct proof explains the positive radial parameters and
the closure at boundary points in the original affine simplex.

For an interior point of $\Delta$, the function
$a/t-b/(1-t)$ decreases continuously from $+\infty$ to $-\infty$.
Choose its root at level $kp$ and put
\[
 S=a/t-kpt=b/(1-t)+kp(1-t).
\]
Multiplying the two equal expressions by $1-t$ and $t$, respectively,
and adding proves $S>0$. Apply the analogous construction to $c,d$
and level $-kq$ to obtain $r$ and $T>0$.
Writing $w=2t-1$, $z=2r-1$, the sum and difference equations are
\[
 1=S+T+kpw-kqz,\qquad
 q=Sw+kpG(w),\qquad p=Tz-kqG(z).
\]
If both caps fail, $S,T>1/2$ and boundedness first gives
$S,T=1/2+O(k)$, $w=2q+O(k)$, $z=2p+O(k)$.
Substitution in the sum improves $S+T$ to $1+O(k^2)$ and therefore
each radial parameter to $1/2+O(k^2)$. The remaining equations yield
\[
 w=2q-kp(1+4q^2)+O(k^2),\qquad
 z=2p+kq(1+4p^2)+O(k^2).
\]
All remainders are uniform on the bounded parameter set. It follows that
$S+T=1+k^2(p^2+q^2+8p^2q^2)+O(k^3)$. The two failed caps require
$S+T>1+k^2(F(p)+F(q))+O(k^3)$, whereas
\[
 F(p)+F(q)-(p^2+q^2+8p^2q^2)
 =4(p^2-q^2)^2+2/1000>0.
\]
This proves coverage of the interior for sufficiently small $k$.
Each sweep is a continuous image of a compact parameter domain and is
therefore closed. Their union contains the closure of the interior, which
is all of $\Delta$. This final compactness step includes the faces and
vertices without extending the fractions $a/t,b/(1-t)$ to endpoints.

\subsection{The barycentric finite enclosure}

Fix such a $k$ and divide $[0,1]$ into $N$ equal closed intervals of
length $\ell$. For a midpoint $t_0$, freeze the first-family direction
and take the perpendicular projection of that interval's compact swept
set as the base. Moving along the frozen direction to the plane $p=0$
gives the $(a,b)$-coordinate intercept
\[
 S(t,1-t)+kp(V(t)-V(t_0)).
\]
With $d=t-t_0$ one has the exact identity
\begin{equation}\label{eq:barycentric-quadratic-identity}
 V(t)-V(t_0)=2d(t,1-t)-d^2(1,-1).
\end{equation}
The intercept is therefore within coordinate distance
$\sqrt2 k\ell^2$ of the sector
\[
 \{\rho(t,1-t):t\in I,\ -2k\ell\le\rho\le M_I+2k\ell\},
 \qquad M_I=\max_I B_k.
\]
This sector is the union of two triangles meeting at the origin. The
absolute Jacobian of $(\rho,t)\mapsto\rho(t,1-t)$ is $|\rho|$;
its area is consequently $\ell M_I^2/2+O_k(\ell^2)$. Its two triangle
perimeters are bounded independently of $I$ and $N$. Thickening a
triangle by a distance $r$ adds at most its perimeter times $r$ plus
$\pi r^2$: cover the added region by the rectangles along its edges
and the circular sectors at its vertices. Apply this separately to the
two triangles with $r=\sqrt2 k\ell^2$. The resulting bound for the
actual compact intercept image is
$\ell M_I^2/2+O_k(\ell^2)$. This enclosure retains the geometric
thickening method; it need not equal an absolute-value polygonal foot.

The plane $p=0$ is parametrized by
$(a,b,(1-a-b)/2,(1-a-b)/2)$. Its Gram determinant is two, so physical
area is $\sqrt2$ times its $(a,b)$ coordinate area. Along a frozen
line parametrized by $p$, the rates are
$ds/dp=kw(t_0)$ and $dq/dp=k(t_0^2+(1-t_0)^2)$.
Using the metric above, its projection cosine is
\[
 [1+k^2 E(t_0)]^{-1/2},\qquad
 E(t)=2w(t)^2+(t^2+(1-t)^2)^2.
\]
There is also a matching lower estimate. Taking $p=0$ in the first
sweep includes the sector
$\{S(t,1-t):t\in I,\ 0\le S\le B_k(t)\}$ in its intercept image.
Its coordinate area is $\tfrac12\int_I B_k(t)^2\,dt$. Therefore the
actual perpendicular base has area at least
\[
 \frac{\sqrt2}{2\sqrt{1+k^2E(t_0)}}
       \int_I B_k(t)^2\,dt.
\]
This lower bound and the preceding geometric upper bound have the same
sum limit. Exchanging the coordinate pairs and changing the sign of the
tilt gives the identical argument for the second family. Thus $2N$
compact-base cylinders cover $\Delta$, and their total area converges
to $2J$, where
\[
 J=\sqrt2\int_0^1
 \frac{B_k(t)^2/2}{\sqrt{1+k^2E(t)}}\,dt,
 \qquad J/\sqrt2=\frac18-\frac{11}{3000}k^2+O(k^4).
\]
For the coefficient, integrate
$F(w(t)/2)/2-E(t)/16$; the two relevant averages are
$\int_0^1 F(w/2)=2/15+1/1000$ and $\int_0^1E=17/15$.
Equivalently the exact change of variable $r=2t-1$ gives
$2J=(\sqrt2/4)Q_R(k)$ with $\beta=1/500$.
Choose $k$ small enough for coverage and $2J<\sqrt2/4$, then take
a finite sufficiently fine partition. The thickened sectors and actual
swept projections need not have equal finite areas; the matching sector
lower bound proves equality of their limiting costs.
